% Bilingual abstract — DRC CoRL 2026: The Validation Gap
% Generated by ars-abstract (academic-paper abstract-only mode, v3.6.7 PATTERN PROTECTION).
% EN: story-style, 215 words. zh-TW: independently composed, 429 characters.
% Do NOT submit the zh-TW to CoRL (English-only venue). Useful for:
%   - Taiwanese / Chinese-language conference versions or journal extensions
%   - Lab internal communication
%   - Broader dissemination of the work

% ─── English Abstract ───────────────────────────────────────────────────────
\section*{Abstract}

Every imitation-learned robot policy reaches the same quiet fork: of the many checkpoints
written during training, which one do you deploy? Measuring success means running the robot,
so the field decides for free, by validation loss on held-out actions --- a habit inherited
from supervised learning, where validation loss and test accuracy rise together. On a robot
they come apart, and the reason is structural. A classifier is graded one example at a time;
a policy is graded by a journey. The small per-step errors validation loss forgives compound,
action after action, until the robot drifts off the states it saw in training and fails.
Validation loss scores the step. Deployment scores the trajectory. We name this the
\emph{validation gap} and study it as a calibration problem --- the way a clinician asks
whether a cheap surrogate tracks the outcome that matters. Pre-registered and reproducible on
consumer GPUs, we score every checkpoint of policies trained across horizons, contact regimes,
and two architectures in two ways at once: offline signals that never touch the robot, and
closed-loop rollouts that measure the truth. One hypothesis organizes the search. If the gap
is born of compounding, then signals of local accuracy should fail just as validation loss
does, while signals of global coherence --- how far a policy drifts from its training
manifold, whether its own rollout stays on course --- should recover what validation loss
misses. The result is a calibrated map: the regimes where a free offline signal can stand in
for a rollout, and the boundary past which you must pay to run the robot. The lesson is easy
to act on. Stop grading robot policies one step at a time. Score the journey --- or, where
you cannot do it offline, run it.

\medskip
\noindent\textbf{Keywords:} validation gap, offline policy selection, imitation learning,
trajectory coherence, robot manipulation, surrogate endpoint calibration

% ─── 繁體中文摘要 (zh-TW) ────────────────────────────────────────────────────
\section*{摘要（繁體中文）}

\begin{CJK}{UTF8}{bkai}

模仿學習的機器人策略訓練會儲存大量檢查點，部署時卻只能選取其中一個。由於直接量測成功率代價高昂，業界慣以驗證損失作為廉價替代——此慣例承繼自監督學習，在該領域兩者高度同步。然而在機器人操作中兩者往往背道而馳，原因根植於評估邏輯的差異：分類器逐樣本評分，策略則以完整軌跡論成敗。驗證損失所寬容的微小單步誤差，在連續動作中逐步累積，最終使機器人偏離訓練分布而失敗。

我們將此現象命名為「驗證差距」，以臨床代理終點校準的框架加以研究：廉價的替代指標，能否忠實追蹤真正重要的結果？本研究採預先登錄設計，跨操作時域、接觸類型與兩種架構，同步為每個檢查點計算離線指標與閉迴路展示的真實成功率。核心假說為：衡量局部動作精準度的指標將如驗證損失一般失效，而衡量全局軌跡一致性者——即策略偏離訓練流形的程度、開迴路重播的穩定性——應能挽回驗證損失所遺漏的預測信號。最終產出一份校準圖譜，指示離線指標可代替展示評估的操作情境，以及何時必須付出代價直接運行機器人。

\medskip
\noindent\textbf{關鍵詞：}驗證差距、離線策略選取、模仿學習、軌跡一致性、機器人操作、代理終點校準

\end{CJK}
